\documentclass[11pt]{article}
\usepackage[T1]{fontenc}
\usepackage[utf8]{inputenc}
\usepackage[margin=1in]{geometry}
\usepackage{enumitem}
\usepackage{hyperref}
\hypersetup{colorlinks=true,linkcolor=black,urlcolor=blue}

\title{Toward Accountable Intuition:\\
A Traceable Live-Play Architecture for Bounded Machine Action on the Intellivision}
\author{Working draft}
\date{October 2026}

\begin{document}
\maketitle

\begin{abstract}
This paper proposes an architecture for bounded machine action in original
videogame play. An Intellivision Master Component remains the authoritative game
and display machine. A separate 6x09-family computer, a future wire-wrapped
HD6309 companion, and a host system observe play, preserve an explicit decision
trace, and only later offer fail-neutral controller actions. We call the first
target an \emph{intuitive program}: a bounded program that identifies
task-relevant differences, ranks permitted candidate actions, and revises that
tendency from observed consequences. This draft does not claim general
intelligence or completed learning results.
\end{abstract}

\section{Introduction}

Machine-learning systems can produce competent play while leaving the route from
observation to action difficult to inspect. Historical computers can make each
bit visible while never adapting their behaviour. This project asks whether an
original videogame platform can support a bounded experiment in which a program
selects what to attend to while leaving enough evidence for another person to
reconstruct why an action was permitted.

The target game system is an original Intellivision Master Component, model
2609. It runs original game software and remains the source of displayed video.
The companion observes the game externally and retains structured history. It
does not begin as a second master of the console bus.

\section{Literature survey and claimed novelty}

The closest technical literature is not a single field but an intersection of
programmatic reinforcement learning (RL), explainable RL, and formalized action
constraints. Programmatic RL represents policies as executable programs rather
than solely as opaque neural networks. Shabadi, Fijalkow, and Matricon study
programmatic policies for gridworlds, including policy-size bounds and a
synthesis procedure \cite{shabadi2024}. Liu et al. use an LLM-guided search to
synthesize programmatic policies in Karel-like tasks \cite{liu2024}. These are
important precedents for the proposal's use of bounded maze policies, but they
do not make a physical action boundary or live hardware provenance their
central object.

Recent work also demonstrates that readable programmatic policies can be useful
under constrained computation. Hu, Zhang, and Baier's ProRL represents
scheduling strategies in a domain-specific language, searches a program space,
and learns parameters while retaining editable policy structure
\cite{hu2026}. This is especially close to the present proposal's interest in
small, revisable policies. It differs in task and substrate: scheduling is not
live play against an original physical game machine, and the work does not
require a human/machine authority gate at an electrical boundary.

Explainable RL provides the second neighbouring field. Saulières surveys the
area through the questions of what is explained and how it is explained
\cite{saulieres2025}. ASQ-IT provides interactive video-based explanations of
RL behaviour using temporal-logic queries \cite{amitai2025}. Such work is
valuable for asking questions of an agent after or alongside its operation. The
present proposal is narrower and more infrastructural: it attempts to make the
observation, candidate set, authorization, act, verification, and fault state
native records of the action system itself, rather than explanations generated
solely after a policy has acted.

The claimed novelty is therefore not that a program can learn a policy, nor
that an RL policy can be explained. It is an \emph{embodied, provenance-first
architecture for intuitive programs}. A candidate-ranking policy is separated
from authorization to act; a physical, fault-neutral boundary and immediate
human takeover constrain later controller emulation; and a replayable trace
spans video observation, finite-state context, candidate ranking, authority,
bounded action, verification, and outcome. The claim remains a design claim
until the protocol yields live trials. It should be evaluated against the cited
fields rather than asserted as an unexplored category.

\section{Intuitive programs}

An \emph{intuitive program} identifies task-relevant differences, ranks a small
set of permitted next actions, and revises that ranking from observed
consequences. It is not thereby general, conscious, or entitled to unrestricted
control. For a maze-oriented task, the intended loop is:

\begin{center}
\texttt{frame/history $\rightarrow$ relevant differences $\rightarrow$ candidates
$\rightarrow$ ranked choice $\rightarrow$ bounded action $\rightarrow$ consequence
$\rightarrow$ revised tendency.}
\end{center}

The first policies may be transparent: wall following, target seeking, or
bounded exploration. A learned host policy may later rank the same candidates.
A ranking is evidence for a choice, not authority to change hardware state.

\section{Accountable action}

The design makes five commitments:

\begin{enumerate}[leftmargin=1.5em]
\item A difference becomes a signal only when it can alter a permitted decision.
\item Comparison is evidence, not an implicit state write.
\item A candidate is not a physical action until authority and guard conditions
      have passed.
\item An act is not established until it has been verified and traced.
\item Ambiguity, timeout, contradiction, and low confidence produce a named
      pause, neutral, or human-takeover state.
\end{enumerate}

The trace is an Ariadne thread: an ordered account of state, candidate,
selection, action, result, error, or halt. It makes a finite controller open to
inspection without confusing that controller with proof of general intelligence.

\section{Apparatus}

The Intellivision 2609 runs the original game and must remain functional when
all companion components are absent, unpowered, reset, or faulted. The M6x09-I
SBC is the operational 6x09 reference platform for RAM-loaded programs, serial
traces, timing experiments, and passive daughterboard work; its RS-232 monitor
path is verified at 19,200 8N1. The wire-wrapped HD6309 board is the primary
future artifact and transfers only demonstrated behaviour from the reference
platform. A host computer captures video, archives trials, and optionally
evaluates candidate-ranking policies. A human player retains immediate takeover
authority.

The M6x09-II SBC remains a useful diagnostic/design reference, but its ACIA
serial acceptance is unresolved and it does not gate this experimental path.

\section{Decision and authority contract}

Each machine turn is a transaction:

\begin{center}
\texttt{target $\rightarrow$ candidate $\rightarrow$ compare $\rightarrow$ guard
$\rightarrow$ bounded action $\rightarrow$ verify $\rightarrow$ next state or halt.}
\end{center}

State, timing, fault, and human/machine ownership guards decide whether a
candidate is lawful. A physical selector exposes Human, Neutral, and Machine
states. Reset, watchdog timeout, host loss, or fault must return a future
controller boundary to disabled/neutral. The first experiments contain no
controller injection: they use external video capture and timestamped human
input. Controller emulation is later work following direct electrical
measurement.

\section{Trace, evidence, and agent guidance}

Each trial records frame identifiers or extracted features; current state and
permitted candidates; candidate ranking and selected action; policy and
firmware versions; authority source; duration; outcome; and error, timeout,
neutral, or takeover transitions. Score is not the sole measure of success.

An agent serves as an engineering shadow, not as a source of hardware truth. It
may organize evidence, cross-check sources, state hypotheses, and propose the
smallest discriminating test. Each material step preserves the observation and
conditions, alternatives, human decision, next test, and correction. This
supports a later podcast account without allowing an agent explanation to become
untested electrical evidence.

\section{Path to empirical results}

Before this manuscript reports performance or adaptation, the project requires:

\begin{enumerate}[leftmargin=1.5em]
\item a preserved M6x09-I S-record load and RAM-execution transcript;
\item a fixed maze-oriented game, configuration, start state, and action set;
\item a transparent baseline policy;
\item replayable trials using the stated trace format;
\item a human or fixed-policy comparison; and
\item evidence that candidate rankings changed after outcomes.
\end{enumerate}

Until then, the contribution is an architecture and protocol for testing
intuitive programs honestly. Once these materials exist, this draft can become
an arXiv-ready empirical report with reproducible artifacts, figures, related
work, and an explicit claim audit.

\section{Conclusion}

The intended threshold is modest but substantive: a machine begins to select
what to attend to and what to try from experience, while another person can
still inspect how and why.

\section*{Contextual materials}
\begin{itemize}[leftmargin=1.5em]
\item The Last Cyberneticist, \href{https://www.spreaker.com/episode/a-machine-must-know-its-next-move--75639925}{``A Machine Must Know Its Next Move,''} 2026.
\item The Last Cyberneticist, \href{https://www.spreaker.com/episode/from-algebra-to-algorithm-multenions-and-program-synthesis--75480441}{``From Algebra to Algorithm: Multenions and Program Synthesis,''} 2026.
\end{itemize}

\bibliographystyle{plain}
\bibliography{references}

\end{document}
